\scanpage{39}
% archive foliation: 38
est déterminée par le complexe de
De Rham-Sullivan de $X_*$, et dans
le cas $k$ quelconque (et cas $2^\circ$ disons
$k$ principal en \unsure{le|cas} $1^\circ$) il est
déterminé par le complexe de
De Rham à p.d. Enfin,
\deleted{il y a} (il est clair que $C^*$
avec les structures suppl. envisagées
\deleted{\unsure{donne}}\inserted{reconstitue} le complexe de De Rham-
Sullivan resp. le complexe de
De Rham à p.d.) Enfin, il
faut voir dans quelle mesure $C^*$
avec sa structure multiplicative
à p.d. permet de reconstituer le
type d'homotopie du $X_*$ dont
il provient. (Enfin, bien sûr,
\unsure{fainéantise} et \unsure{platonisme|platitudes} --- !!!)
